\documentclass[11pt]{article}
\usepackage[T1]{fontenc}
\usepackage[utf8]{inputenc}
\usepackage[margin=1in]{geometry}
\usepackage{lmodern,microtype,amsmath,amssymb,enumitem,xurl}
\usepackage[hidelinks]{hyperref}
\hypersetup{pdftitle={Operator Archaeology: Procedural Media, Reasoning, and Authority},pdfauthor={Brian Cameron}}
\setlength{\emergencystretch}{2em}
\setlist{nosep}
\title{Operator Archaeology:\\Procedural Media, Reasoning, and Authority\\
\large From Historical Practices to AI}
\author{Brian Cameron}
\date{September 2026}
\begin{document}
\maketitle
\begin{abstract}
A board, calendar, diagram, or table can preserve more than an image or a belief.
It can preserve a way of doing something. This essay introduces operator
archaeology as a protocol for investigating such procedures while distinguishing
historical evidence from modern models. Selected examples from astronomy,
games, Egyptian funerary texts, Ramon Llull, cryptography, and John Dee lead
to a contemporary question: when a system organizes reasoning or produces
an answer, what supports treating that answer as knowledge or as a basis for
action? The historical cases have different purposes and evidence; they are
not stages of one hidden tradition. Their comparison makes three distinctions
clearer: procedure and meaning, representation and its source, and generated
output and warranted authority. These distinctions connect the author's
historical research with a separate program on reconstructible reasoning,
bounded correction, and accountable AI use.
\end{abstract}

\section{Reasoning outside the head}

Put a game board on a table. A player moves a piece, and the situation changes.
The board preserves where the pieces are; the rules determine which moves are
available. Some of the work of remembering and reasoning has moved into an
arrangement that the players can inspect together.

Put a timetable beside it. A reader follows a row to a destination and a column
to a departure time. This object preserves a different procedure. Put a
calculator beside both. Its operations are different again, although all three
objects let a person work through a problem with something outside the head.

These ordinary examples suggest a question that becomes more difficult, and
more interesting, when the users and their instructions are no longer present.
What work did a surviving structure allow someone to perform? Which parts of
that work remain recoverable? Which parts depended on a teacher, convention,
institution, or practice that the surviving object does not preserve?

I call the investigation of these questions \emph{operator archaeology}.
An operator, in this essay, is a specified transformation or rule of use:
selecting positions, rearranging letters, advancing a numerical scheme,
permitting a move, or organizing a sequence of actions. Sometimes a historical
source states the procedure. Sometimes the term names a modern hypothesis
that still needs testing. The difference must remain visible.

My interest comes partly from software architecture and application security.
A system can execute its instructions perfectly and still produce an
inappropriate result because its assumptions, inputs, or permissions were
wrong. That habit of separating questions is useful here. It encourages me to
ask what an artifact supports, what my representation adds, and which claims
depend on evidence I do not have. It does not give software concepts historical
authority over the sources.

The larger question is about instruments of thought. How do people externalize
reasoning, and how do those instruments come to guide judgment? This question
has a history, but it also has immediate force when a language model supplies
an explanation, a risk assessment, or a proposed action.

The connection is a comparison of problems. A wheel with explicit terms, a
learned statistical model, and a ritual instruction can all organize inquiry
while doing so through very different means. Their differences help us ask
which assumptions belong to the instrument, which belong to the user, and
which belong to the setting in which an answer is accepted.

There is established work behind this approach. Edwin Hutchins's account of
navigation moves the unit of analysis beyond the isolated individual to a
team, its instruments, and its practices.\cite{hutchins} Source criticism and
reconstruction scholarship likewise require attention to mediation and
uncertainty. The London Charter, for example, asks cultural-heritage
visualisations to distinguish evidence-based restoration from hypothetical
reconstruction and document interpretive decisions.\cite{london}
Operator archaeology applies comparable disciplines to proposed procedures
across media. Its contribution is an organizing protocol and a set of explicit
questions, rather than a claim to have invented careful historical reasoning.

\section{Ask what a structure does}

A surviving table is an observation. A proposed way to traverse it is an
interpretation. A successful extraction under that traversal is a result.
A claim that its historical makers intended the extracted message is a
further claim. These steps often appear together in an exciting reconstruction,
but they require different kinds of support.

The first task is to identify the source. Which artifact, manuscript,
page, edition, or transcription is being used? Has its layout been flattened?
Are row numbers original or editorial? Are damaged marks being treated as
blanks, zeros, or guessed letters? A model begins making choices before its
first calculation.

Next comes the proposed procedure. What are its inputs? What changes? What
does it preserve? What result would count as success? A claim about a
repeatable transformation should be specific enough for someone else to try
it, including the choices that make an attractive output disappear.

The strongest procedure is not necessarily the most elaborate one. If an
ordinary indexing convention explains a table, a hidden program adds no
explanatory value until it predicts something the indexing account cannot.
If a ritual text specifies a sequence, recovering that sequence need not
require treating every name or number as a computational instruction.

A reconstruction also needs a boundary. A source can attest an instruction
without proving that it was executed, effective, or universally followed.
A modern simulation can show that a proposed procedure is possible without
establishing that anyone used it historically. A useful analogy can suggest
a test without establishing transmission between cultures.

\subsection{Seven habits of a correctable reconstruction}

The detailed protocol develops seven principles.\cite{method} In public-facing
terms, they amount to the following habits:
\begin{enumerate}
\item \textbf{Keep levels separate.} Say whether a claim concerns a source,
an observed feature, a model, a reconstruction, or historical meaning.
\item \textbf{Show what is missing.} Preserve the distinction between a damaged
record, an uncertain reading, and an actual feature of the object.
\item \textbf{Correct locally.} When one proposed operation fails, narrow or
replace that operation rather than reshape the entire theory to rescue it.
\item \textbf{Make the idea testable.} State what a new example should show and
what would count against the proposed rule.
\item \textbf{Compare operations precisely.} Two circles do not establish a
shared function. Specify what is being compared and the evidence on both sides.
\item \textbf{Separate discovery from justification.} A failed explanation may
suggest a useful idea. That idea needs its own argument or evidence.
\item \textbf{Ask whether the procedure fits its setting.} Consider the medium,
skills, resources, institutions, and purposes that would make it usable.
\end{enumerate}

These habits are not a demand to weaken every claim. They allow confidence
to be placed where the evidence earns it. A narrowly demonstrated operation
can be described strongly while a larger interpretation remains open.

They also make uncertainty productive. Instead of saying only that a table
is mysterious, an analyst can identify the unresolved question: whether a
mark specifies a reading direction, whether a boundary resets a schedule,
or whether two witnesses preserve different versions of an instruction.
Those questions tell us what source work or experiment would help.

A practical miniature shows the difference. Suppose a manuscript appears to
instruct the reader to retain alternate words and take their initials.
The analyst must fix where counting starts, what counts as a word, and whether
a line break changes the count. If the desired output requires correcting
three spellings after it has been recognized, those corrections belong in
the result. If an independent example works without them, the rule has gained
support. If every new example requires a different repair, the proposed
generality has weakened.

This miniature is hypothetical. It explains the discipline; it is not evidence
for a particular historical cipher. A source-grounded claim would additionally
identify the manuscript instruction and the carrier being tested.

The method is deliberately cross-media. It asks what work a structure performed
before assigning it to a modern disciplinary category. That makes room for a
procedure carried jointly by text, arrangement, speech, and learned practice.
It also makes room for the conclusion that the best explanation is not an
operator model at all.

\subsection{What survives a failed explanation}

A speculative interpretation can fail without making every question it raised
worthless. Suppose a proposed decoding depends on a reading direction that the
source does not support. The decoding loses its historical warrant. The work
may still suggest a useful way to compare reading directions in another,
independently documented case.

What survives is a candidate idea, not a transferable result. It should enter
the next investigation as a question to test. If the new case supports it,
that support belongs to the new case. The original failure is not retrospectively
converted into a discovery.

This distinction matters to a wide research program. Curiosity generates
connections more quickly than evidence can establish them. A disciplined
program needs somewhere to keep those connections without making the active
argument depend on them. An archive or a speculative companion can preserve
a question while the main paper states what has actually been supported.

That arrangement protects the imaginative part of research as well as its
credibility. A researcher can explore a possibility, record why it failed,
and move on without either concealing the failure or abandoning the whole
subject. Correction becomes a way to refine the inquiry.

\section{Time, play, and passage}

A historical arc does not need every example to do the same thing. It becomes
more useful when different cases show the limits of different abstractions.
Astronomical calculation, board games, and funerary texts offer three ways
to investigate procedure without treating all procedure as computation.

\subsection{A numerical scheme is more than a list}

Lis Brack-Bernsen's study of Babylonian lunar column $\Phi$ discusses a
linear zigzag scheme and its relation to observations and other astronomical
quantities.\cite{brack} This is a strong setting for procedural comparison:
the object of study includes numerical relationships that can be investigated,
rather than a pattern noticed only in an evocative image.

For a modern teaching model, imagine a sequence that increases by a fixed
step until an upper bound and then decreases. The current value, direction,
step, and bounds matter. Merely listing the output values hides the rule
that connects them. This toy model is not a transcription of a Babylonian
column; it illustrates what it means to distinguish stored results from a
specified procedure.

It also illustrates why an adequate explanation needs more than the word
\emph{cycle}. Two sequences can repeat while preserving different information.
One may mark an occasion; another may approximate a changing quantity.
Their outputs can look similar at a coarse scale while their operations and
historical uses differ.

The focused Mesopotamian paper examines these distinctions in greater
detail.\cite{mesopotamia} For the present essay, the important point is
methodological: a numerical table is especially informative when one can
ask how its entries depend on one another, how the quantities were identified,
and where reconstruction remains uncertain.

A successful replay would then answer a bounded question. It could show that
declared rules reproduce specified entries. It would not by itself establish
the complete process through which the historical theory was discovered.
The difference between construction, performance, and discovery survives
even in a comparatively well-documented mathematical case.

This matters for contemporary systems as well. We can reproduce a calculation
without knowing why a parameter was chosen. We can reproduce a model's output
without justifying its use in a new setting. Reproducibility answers an important
question, but it does not answer every question about an instrument.

\subsection{The board and the rules are different evidence}

The British Museum preserves both a game board from Ur and a much later
cuneiform tablet describing the twenty-square game.\cite{urboard,urrules}
The temporal distance matters. A later rule text is valuable evidence,
but it cannot automatically settle every rule used with every earlier board.

The methodological lesson is unusually clear. A board preserves possible
positions and visible markings. A rule text preserves instructions. A modern
playable reconstruction combines these with interpretive choices. Each
contributes something the others do not guarantee.

When rules are fixed, a game offers an intelligible example of state and
transition. A move can be legal or illegal under those rules. The result of
a chance event and the permission to act on it are distinct. A piece's position
does not mean the same thing without the conventions that govern it.

There is no need to infer that the players understood formal state machines.
The modern description can be useful because it makes our reconstruction
inspectable. It tells us what changes when a rule changes, and which conclusions
depend on that choice.

A hypothetical reconstruction might give a marked square protection from
capture. Another might give it an extra turn. Both are playable, and both
could make the board feel purposeful. Playability alone would not choose
between them. A historical instruction, a related witness, or another
independent constraint would be needed.

This is a practical antidote to an attractive error: treating the coherence
of one's own model as proof that it is the historical model. The better
reconstruction exposes its decisions. It lets another reader ask whether
the evidence supports those decisions.

Games can also inspire experiments about learning, strategy, or cooperation.
Such experiments would study particular players under specified conditions.
They would not establish that an ancient game was designed to teach the
modern lesson we happen to find in it. Recovering a useful educational
primitive and recovering an original educational purpose are separate tasks.

\subsection{Knowledge at a threshold}

Rita Lucarelli's study of guardian figures in Book of the Dead spells
144--147 describes encounters at guarded passages and the importance of
knowing names. It also examines adaptation of such material in temple
contexts.\cite{lucarelli} The evidence concerns religious compositions
with their own meanings and histories.

An analyst can nevertheless ask procedural questions. What encounter is
specified? What knowledge or speech is required? Who is addressed? What
does successful passage mean within the composition? These are more informative
questions than identifying every gate with a modern security check.

Here the analogy must stop short of explanation. A contemporary access-control
system and a funerary declaration can both involve conditions for passage.
That relational comparison does not make the religious encounter an early
authentication protocol, nor does it tell us whether the action was performed
as the text depicts.

The religious significance is part of what must be explained. Names, beings,
places, and declarations cannot simply be stripped away to reveal an allegedly
real computational core. An abstraction that removes them may help compare
one relation while becoming inadequate for a question about the composition's
meaning.

These cases therefore support three different lessons. Numerical schemes
make relationships among values explicit. Games make rules and possible moves
inspectable. Passage texts organize encounters in which knowledge and action
have religious significance. Calling them procedural media identifies a
question about their use. It does not give them a common ontology.

\subsection{When a machine really is a machine}

The Antikythera mechanism offers a useful contrast. Freeth and colleagues
investigated surviving gears and inscriptions to reconstruct features of
an ancient Greek astronomical calculator.\cite{freeth} Here mechanical
relations are material evidence, not only a metaphor applied by the analyst.

The contrast disciplines the comparison. Evidence for gears in one object
does not turn a wheel drawn in another text into a mechanical device.
Conversely, the absence of gears does not prevent a written procedure from
being sophisticated. Different media can preserve different parts of a task.

The interesting question is what the medium makes stable. A gear relationship
can constrain motion; a diagram can constrain an arrangement; a text can
state a rule; a trained practitioner can preserve a convention that the object
alone omits. Recovering the relevant combination requires attention to the
particular case.

\section{Llull and the ambition to organize thought}

Ramon Llull is central to this essay because the question of externalizing
reasoning becomes especially explicit in his Art. His philosophical and
religious project sought to establish truth through an organized method of
combination. That project included Christian missionary purposes; it was
not a neutral catalogue of possible ideas.\cite{priani}

In the \emph{Ars brevis}, the Fourth Figure has three circles, with the outer
one fixed and the inner two movable. Letters can be brought into combinations.
The text also requires the practitioner to consider their meanings and
their applicability to the subject.\cite{llull}
The arrangement is therefore only part of the procedure. A person still
has to understand and use the terms.

The attraction is easy to appreciate. A difficult subject can overwhelm us
because relevant ideas arrive in an accidental order. Some are remembered,
some are forgotten, and some are repeatedly paired while others are never
considered together. An organized combinatorial device offers a way to
make relations available for inspection.

To see the analytical point, imagine a modern exercise using three concepts:
trust, competence, and accountability. Their arrangement might prompt questions
about competent but unaccountable action, accountable but unreliable advice,
or trust earned through demonstrated competence. This is my teaching example,
not a reconstruction of Llull's alphabet or a demonstration of his Art.
It shows how a constrained arrangement can prompt inquiry without already
containing an answer.

The distinction is between making a question available and establishing a
conclusion. A procedure that reliably produces combinations has accomplished
something real. It may enlarge the space a reader examines or make omissions
easier to notice. But the argument connecting a combination to a conclusion
requires additional work.

That distinction is a contemporary analytical lesson from the case. It should
not be mistaken for a description of Llull's own standard of proof. A historical
account must explain his definitions, rules, and theological commitments on
their own terms. The public comparison asks what becomes visible when we
distinguish those commitments from the mechanism of combination.

\subsection{The vocabulary is part of the instrument}

A reasoning device begins before its first operation. Someone decides which
terms it can represent. Someone decides whether different uses of a word
belong together, whether a question is admissible, and which relations matter.
Those choices shape what can be asked through the device.

This is why a combinatorial system cannot be evaluated only by counting the
combinations it produces. An exhaustive arrangement over a narrow vocabulary
can still omit the issue that matters most. A beautifully organized set of
premises can make a disagreement difficult to express if the disagreement
concerns the premises themselves.

Imagine an instrument for evaluating a public proposal that represents only
cost and speed. It could enumerate their relations perfectly. Yet questions
about consent, distribution of benefits, or irreparable harm would enter only
if someone supplied a way to represent them. The shortcoming would precede
any error in combination.

This is a hypothetical example of representation adequacy. It does not accuse
Llull of the modern omission. It uses the contrast between vocabulary and
operation to identify a question that belongs to any instrument of reasoning:
what does it allow us to distinguish?

The more general an instrument claims to be, the more important that question
becomes. A local procedure can announce its restricted purpose. A purportedly
universal one must explain why its concepts and rules suffice across the
domains it claims to address.

This issue connects Llull to AI without calling him an early language-model
designer. The relationship concerns how a system organizes an inquiry and
what commitments it imports. The technologies, purposes, and sources of
their representations remain different.

\subsection{A productive comparison with AI}

When an AI system generates several explanations, those explanations become
available for evaluation. Their availability is useful. Their fluency,
agreement, or mutual consistency does not establish that they are supported
by independent evidence.

A combinatorial instrument helps make one version of that distinction easy
to see. Its rules may explain why a candidate appears. That account of
production is not yet an account of why the candidate should be accepted.
The same separation is necessary when the production process is much more
complex and less transparent.

We should therefore ask two questions of an instrument that organizes thought.
Does it help us generate or inspect relevant possibilities? Does it help us
evaluate them against evidence and appropriate standards? An instrument may
be excellent at the first and limited at the second.

It can still be worthwhile. Suggesting a neglected question, exposing a
missing relation, or giving a discussion a manageable form are significant
contributions. The danger appears when those contributions silently become
a claim that the instrument has settled what should be believed or done.

Llull's case belongs in this essay because it keeps the ambition visible:
to turn reasoning into something organized, teachable, and repeatable.
A careful historical account can preserve that ambition while a contemporary
analysis asks where production ends and judgment begins.


\section{Trithemius and Dee: procedure, record, interpretation}

Early modern cryptography brings a different advantage. In some cases an
extraction can recover an independently checkable text. This gives the
analyst a more discriminating target than an interpretation that remains
plausible regardless of the output.

Jim Reeds's study of Book III of Trithemius's \emph{Steganographia} presents
numerical cryptograms and recovered plaintexts within material presented
in an angelic and astronomical idiom.\cite{reeds}
The result supports specific cryptographic procedures. It does not establish
that every number, sign, or ritual feature elsewhere in the work has the
same role.

This case sharpens the meaning of doing something. A decoder takes
specified input through a declared transformation and produces a result.
Whether the work also carried religious significance, how its readers
understood it, and which procedures they used are additional historical
questions. A technical result need not be made to answer all of them.

It also sharpens the test. A striking phrase is insufficient if the analyst
had unlimited freedom to change selection, spelling, phase, or reading direction.
The persuasive reconstruction makes the procedure explicit and shows why its
output is supported beyond the analyst's taste.

The focused Steganographia paper separates source material, declared
transformations, replayed fixtures, and unresolved components.\cite{trithemiusPaper}
That separation is valuable even before a complete reconstruction exists.
It tells a future investigator where another witness, clearer instruction,
or independent carrier would change the evidentiary position.

\subsection{A connection is not a transferred mechanism}

The National Library of Wales identifies a 1591 transcript of
\emph{Steganographia} attributed to Dee's hand and describes his earlier
pursuit of a copy in 1563. It leaves his understanding or use of the
particular system uncertain.\cite{nlw}
This is a documentary connection, with limits.

Possession makes borrowing possible. It does not identify a borrowed rule.
The later surviving copy cannot by itself establish which procedure Dee
used in an earlier conversation or table. Shared spirits, letter arrays,
or circular arrangements supply questions for comparison, not a completed
genealogy.

This suggests a useful maxim: results may travel without their generating
procedures. A copied diagram might preserve an output without its explanation.
It might instead preserve enough instruction for a knowledgeable reader to
reconstruct a practice. It might be adapted to a new purpose. These are
separate possibilities to investigate.

A useful transmission claim identifies what moved. Was it wording, an
arrangement, a method, a trained practice, or an interpretation? The evidence
needed to answer varies. Similarity at the level of a printed shape is
too coarse to settle all of those questions.

\subsection{A local case from Dee's record}

Dee's record for June 18--19, 1583 provides a more bounded example than a
proposed universal angelic machine. In Casaubon's edition, reported perceptions,
Dee's questions, and replies attributed to Galvah are distinguished.
The exchange specifies a table arrangement while leaving colors to Dee's
discretion; the next day's entry records acceptance of his placement in
small triangles.\cite{deeLayout}

This permits a modest but worthwhile analysis. Prescription, discretion,
inscription, and reported acceptance are different parts of the recorded
activity. The record makes some of those distinctions visible. It does not
independently establish the source of the attributed reply.

The local discretion about colors also constrains a tempting hypothesis.
It supplies no basis for treating every color as a cipher selector.
Another artifact might give a color a specific procedural role, but that
would require its own instruction or evidence.

In a separate Heptarchic example, a source describes removing a common
initial letter, arranging paired names, and reading according to an
orientation.\cite{heptarchia}
A modern description can distinguish removal, indexing, and permutation.
That describes a local rearrangement. It does not decode the theological
meaning or establish a generator for another collection.

These cases make the investigation more interesting, not less. They replace
a broad question about a hidden machine with tractable questions about how
specific materials were constructed, recorded, and interpreted.

\subsection{The record is itself an instrument}

A written account preserves an exchange through choices about wording,
speaker labels, arrangement, and selection. It can later become the basis
of a table, interpretation, or dispute. The record is therefore part of the
history to investigate, not a transparent window through which every stage
of an event can be independently inspected.

For a bounded source study, we can ask who is represented as questioning,
reporting, writing, or correcting, and how those roles change in the relevant
entries. We can compare copies and editorial interventions where available.
We can investigate whether an instruction specifies a reproducible operation
or whether the operation remains underdetermined.

Those questions do not require a diagnosis of the participants or a verdict
on their experiences. Nor does distinguishing production stages establish
that the stages caused an effect on belief. A claim about perceived agency
or credibility needs evidence appropriate to that claim.

The useful general distinction is between producing an inscription, interpreting
its content, and attributing its source. One can describe a known arrangement
without explaining all three. A future speculative model might propose a
mechanism for one stage, but its success would need to be demonstrated
rather than borrowed from the existence of the others.

That is the stance taken in the focused Dee paper.\cite{deePaper}
It preserves religious context and source-attested procedure while leaving
larger mechanisms open. The historical case contributes to the public essay
precisely because those limits remain visible.

\section{Leibniz and the difference between a calculus and a model}

Leibniz's project of a universal characteristic and a logical calculus
provides a further historical point of comparison: an ambition to express
reasoning through symbols and rule-governed operations.\cite{leibniz}
This essay does not reconstruct that project in detail. Its relevance is
the explicit relationship among representation, calculation, and judgment.

Imagine that a disagreement could be translated into symbols, with agreed
rules for operating on them. A calculation might then help determine whether
a conclusion follows. But agreement about the symbols and their interpretation
would matter. So would the relation between the represented propositions
and the world to which the conclusion is applied.

A procedure can be correct relative to its premises while its application
is mistaken. It can establish a conditional conclusion without establishing
the premises. And a conclusion about what is the case need not establish
who is entitled to act.

These are philosophical distinctions, not a claim that Leibniz solved
or ignored every contemporary version of them. They explain why the ambition
to make reasoning calculable remains interesting even when the instruments
change.

Current language models should not be described as Llullian wheels or as
an implementation of Leibniz's universal calculus. The Transformer paper
introduced an attention-based architecture for sequence transduction;
subsequent instruction-following research includes training with human
feedback.\cite{vaswani,ouyang}
These are specific technical developments with their own evidence.
The historical comparison does not explain their training or operation.

A language model can nevertheless enter the same practical situation as
another instrument of thought: someone supplies a question, receives an
output, and decides what weight to give it. The instrument's production
process is one part of the situation. Evidence, interpretation, permission,
and consequences are others.

The continuity is therefore at the level of questions we must ask.
What has been represented? What operation produced this answer? What supports
its application? How can error be discovered and corrected? Which decisions
remain outside the instrument's competence or authority?

That comparison lets historical examples illuminate AI use while retaining
the differences that make each example worth studying.

\subsection{What changes when answers become cheap}

A modern language model can make trying an explanation relatively easy.
A user can ask for alternatives, objections, analogies, or different styles
of presentation in the same interaction. This is an opportunity to organize
inquiry, provided the output's role is specified.

Consider a hypothetical research workflow. An assistant suggests five
interpretations of a table. One interpretation prompts the researcher to
find a source they had overlooked. Another generates a test that rejects
the preferred rule. A third is eloquent but has no evidentiary basis.
The same interaction has produced useful assistance and unsupported content.

Treating all five as findings would discard the distinction that makes the
assistance worthwhile. The researcher needs to keep track of what the model
proposed, what the source supplied, what was tested, and what remained an
interpretation. The ease of generating another paragraph does not reduce
that need.

This suggests a productive role for AI in historical inquiry. It can help
articulate alternatives, organize source questions, compare declared models,
or prepare a reproducible experiment. The warrant for a claim still depends
on the sources and the analysis. An assistant's account of why it produced
an idea is itself another output to examine, rather than privileged access
to its internal causes.

The point extends beyond historical research. An instrument that makes
answers easy to obtain can also make the distinction between proposal and
decision easy to overlook. The workflow must preserve that distinction
rather than assume that users will recover it after the answer has been
acted upon.

\section{Representation, validation, and authority}

The most important contemporary lesson concerns a failure that can occur
before a system makes an inference. It can choose a representation that
removes the distinctions needed for the decision, then reason correctly
over what remains.

Consider a hypothetical inspection record with three states:
\emph{compliant}, \emph{noncompliant}, and \emph{not assessed}.
A dashboard displays only pass and fail. If not assessed is folded
into either category and that display is used to authorize an action,
a consequential distinction has disappeared.

The loss is not repaired by perfect logging. A hash-bound record can preserve
the binary classification faithfully. A deterministic workflow can apply
the gate exactly as specified. A successful audit of those transitions can
still leave the original evidentiary problem unresolved.

The representation must therefore be evaluated in relation to the proposed
use. Does it preserve the distinctions material to that decision? A simplified
display may be adequate for counting completed inspections and inadequate
for declaring every unflagged item safe. The same abstraction can be
appropriate for one query and inappropriate for another.

This is \emph{representation adequacy}, a principle developed in my reasoning
papers.\cite{reliability,rprp} It is not a demand to preserve everything.
Abstraction would be impossible under that demand. It is a demand to identify
which losses matter before promoting a derived view into a stronger claim.

\subsection{A view does not inherit all the properties of its source}

A summary, category, score, or extracted token can stand for part of an
object without becoming interchangeable with it. A map can support a route
while omitting features that matter for another journey. A total can support
an accounting query while concealing the distribution needed for a fairness
question.

The problem is a silent promotion. Something designed as a convenient view
is treated as if it retained the evidentiary or semantic force of its source.
The analyst then reasons over the promoted representation without revisiting
the missing distinction.

Return to the historical examples. A linear transcription can support a
question about word order while losing the placement of a marginal instruction.
A modern model of a game can support simulation while leaving the historical
rule choice unsettled. A procedural abstraction of a religious composition
can preserve a relation while omitting what matters for its interpretation.

The common principle does not depend on any of those historical claims.
Each application needs its own evidence. The comparison helps us recognize
the form of the problem and ask the next question more precisely.

For AI use, an explicit practice would identify the source material, the
derived representation, the intended decision, and consequential losses.
If the necessary distinction is unavailable, the answer may need to remain
conditional or the transition may need to wait. The appropriate response
depends on the use, its consequences, and the possibility of obtaining
better information.

This assessment is itself fallible. It can miss a distinction, misunderstand
the intended decision, or overestimate the possibility of recovering lost
information. Calling a representation adequate should therefore record
the use for which it was assessed and the grounds for that assessment.
It should not become an unrestricted certificate for every later use.

\subsection{Production is not justification}

An answer can be insightful because it connects ideas, proposes a useful
distinction, or exposes a neglected possibility. It can be valuable before
it is established as a historical or factual claim.

That does not eliminate the need for justification. An idea's origin and
the reasons for accepting it are different. A failed reconstruction might
inspire a good concept. A mistaken explanation might point toward a testable
question. The resulting concept or question must earn support on its own.

The same discipline applies to agreement among generated answers.
Three explanations from one system do not automatically become three
independent witnesses. Different roles or prompts can be useful for criticism,
but the independence and evidentiary basis of their claims must be examined.

A reader should be able to distinguish the claim, its supporting material,
the inference connecting them, and the uncertainty that remains.
That makes a disagreement actionable. Another reviewer can challenge a
reading, source, assumption, or inference without having to reject the
entire research program.

This is why correction should be local where possible. A model that explains
one artifact may fail on another. A particular selector may fail while the
broader question remains worthwhile. A source-grounded rearrangement can
survive even when a proposed unified mechanism does not.

The discipline can strengthen confidence. Once dependencies are explicit,
a strong claim need not carry every speculative extension around it.
Its support can be assessed at the appropriate scale.

\subsection{Accuracy does not grant permission}

Suppose an AI accurately estimates a consequence. That accomplishment does
not by itself authorize it to choose the action. The proposed action may
require consent, a mandate, professional judgment, or an institutional process.

Conversely, an authorized actor can make an inadequately supported decision.
Permission and information are different requirements. Neither should
be inferred from the other.

This matters because an instrument can acquire practical influence through
repeated reliance. A recommendation becomes the default; departing from it
requires effort; eventually the workflow treats its output as decisive.
At that point the questions about authority concern the surrounding arrangement
as well as the model's text.

This is a hypothetical organizational pattern, not a measured finding about
every deployment. Its value is to identify what a concrete audit would examine:
who can contest an output, what reasons are required, where responsibility
sits, and whether an exception is practically available.

My governed-bundle work treats these as design questions.\cite{bundles}
An execution record can help establish what happened. It cannot alone
establish the adequacy of the representation, the quality of the judgment,
or the legitimacy of the action. Those dependencies must be supplied
rather than inferred from successful execution.

An instrument's role also needs to be distinguished from its intended role.
A designer may intend to provide suggestions while a surrounding organization
uses them as binding classifications. Observing that use does not establish
that the instrument was built for it or that the use is legitimate.
Function, intention, and authority need separate accounts.

\subsection{Correction, reversal, and repair}

An explanation can be corrected after its use. A decision may be reversible.
Its consequences may still need repair. These are different capacities.

A corrected assessment does not automatically restore an opportunity lost
through the original assessment. Reversing a state transition does not
necessarily reverse what people did in response. A log can preserve an
account of the error while the means to repair it remain absent.

That distinction gives bounded correction a practical limit.
A workflow should ask not only whether its record can be revised, but what
can still be changed and who can address residual consequences. Some choices
can wait for better evidence; others have deadlines; some consequences are
irreversible.

No numerical threshold answers all of these questions. Judgments about
acceptable loss, consent, and repair involve normative premises and affected
parties. A dependency graph can clarify where a decision rests without
determining the values that should govern it.

The philosophical papers \emph{Between Caves} and \emph{The Controller}
develop this separation between informational constraints, judgment,
and responsibility.\cite{caves,controller}
Their arguments stand on their own; the historical comparisons here provide
illustrations and questions, not the foundations of those arguments.

\section{What a thinking instrument must make visible}

The historical cases offer a way to keep the instrument in view.
A board requires rules and players. A table requires identification of its
quantities and a convention of use. Llull's figures require meanings,
definitions, and practice. A cryptogram requires a particular transformation.
A recorded exchange requires attention to its mediation.

For an AI-assisted workflow, the analogous task is to identify the arrangement
that turns an output into a decision. What material entered the system?
How was it represented? What did the model contribute? Which claims were
checked independently? Who supplied the authority to act? What happens
when the result is challenged?

These questions can be made concrete without promising a universal checklist.
For a proposed historical reconstruction, preserve the witness and the
analyst's choices. For a classification system, identify the distinctions
its labels retain and discard. For a recommendation, identify the warrant
for treating it as advice and the separate basis for acting on it.

There is also a question about practical fit. A documented procedure can
depend on training, equipment, institutional support, or a convention that
a new user lacks. Copying its instructions does not automatically recreate
its conditions of use.

A modern workflow might demand independent review while providing no time
or access for it. It might require evidence that a participant cannot obtain.
It might assign responsibility to someone who cannot stop the transition.
Those are hypothetical design failures to investigate, rather than problems
solved merely by writing the requirement.

The artifact is therefore not always the whole instrument. To understand
what it can do, we may need to reconstruct the relation among media,
people, skills, and institutions. That is a research task in history and
a design task in contemporary systems.

A useful instrument makes its contribution specific. It can enumerate
possibilities, preserve a state, calculate a value, propose an explanation,
or organize a discussion. Its limitations need to be expressed at the
same level of specificity. A general instruction to use judgment is too vague
when the workflow does not say which judgment is required or who can exercise it.

The central risk is not that an external instrument contributes to thinking.
Such contributions can be valuable. The risk is that the contribution becomes
difficult to distinguish from the grounds for accepting the result.
Good design keeps that distinction available to the people who rely on it.

\subsection{A modest workflow}

Consider a hypothetical assistant helping review a proposal. The assistant
first identifies the actual question: whether the proposal satisfies specified
requirements, rather than whether its language sounds reassuring. It records
which requirements and evidence it has. It then produces an assessment whose
uncertainty remains attached to the particular claims.

A separate review can challenge the source choices, the interpretation of a
requirement, or the representation that supports the conclusion. If a source
changes, the affected claims can be reconsidered. If one requirement is not
assessed, the resulting view preserves that status instead of quietly treating
it as satisfied.

Only then does an authorized participant decide what the assessment permits.
An unresolved question might block an action, require a narrower action,
or be accepted under a declared exception. The assessment does not invent
that authority. The decision records the grounds and responsibility for it.

This is a proposed design pattern. It has not been evaluated as an intervention
in this essay. Its purpose is to show how the distinctions can shape a workflow
rather than remain philosophical slogans. Its performance would depend on
the actual requirements, sources, reviewers, and institutions using it.

The arrangement also gives AI a useful role. The model can help find overlooked
dependencies, articulate objections, or prepare a comparison. Those
contributions need not be diminished by refusing to treat them as self-validating.
They become easier to use when their place in the process is clear.

\section{What the method earns}

Operator archaeology earns its place when it produces a better question,
a clearer reconstruction, or a more discriminating test. It does not earn
its place by redescribing every artifact in the same vocabulary.

The strongest historical claims in this essay remain local. A particular
rule text bears on a game. A studied numerical scheme bears on an astronomical
quantity. Llull's figures and instructions bear on his Art. Reeds's
cryptanalysis bears on specified cryptograms. Dee's recorded instructions
bear on specific arrangements.

Their comparison is useful because the differences remain intact.
They let us distinguish storing a result from preserving a procedure,
arranging concepts from establishing an argument, and recording an attributed
message from validating its source.

A broader story can still be told. People repeatedly create media through
which memory, calculation, interpretation, and action become organized
outside an individual mind. Sometimes the procedure is explicit; sometimes
part of it resides in practice; sometimes our reconstruction remains uncertain.
The story has several paths and purposes. It need not be forced into a
single developmental sequence.

That is also why this essay keeps its historical and contemporary arguments
connected but separable. A philosophical principle should not depend on
an unverified claim about an ancient object. A historical reconstruction
should not gain credibility merely because it resembles a useful modern
design.

The detailed history stack supplies evidence owners and narrower questions.
The method paper supplies the reconstruction protocol. The reasoning papers
develop conditions for representing, checking, correcting, and governing
contemporary reasoning. This essay is an entry point into those projects,
not a substitute for their evidence or arguments.

Its central proposal is simple enough to carry between them:
ask what a structure enables someone to do, preserve the distinctions
needed to understand that work, and identify the grounds on which its
outputs may guide belief or action.

An instrument can help us think without settling every question for us.
The task is to make that help inspectable, its scope explicit, and its
consequences open to responsible judgment.

\section*{Further reading}
The companion \href{https://github.com/yippibrian/history}{history repository}
contains the focused historical papers and their source-specific limits.
The separate \href{https://github.com/yippibrian/reconstructibility-under-projection}{reasoning repository}
contains the technical and philosophical work on reconstructibility,
representation, and governed use. These links provide routes into the
research; external sources supporting the examples are listed below.
No new historical replay, corpus experiment, or user study is reported
in this essay.


\begin{thebibliography}{99}
\bibitem{hutchins}
Edwin Hutchins. \emph{Cognition in the Wild}. MIT Press, 1995.
Author's introduction and chapter summaries:
\url{https://pages.ucsd.edu/~ehutchins/citw.html}.
\bibitem{london}
The London Charter. \emph{The London Charter for the Computer-Based
Visualisation of Cultural Heritage}, version 2.1, 2009.
Principle 4, Documentation, especially 4.4--4.10.
\url{https://londoncharter.org/principles/documentation.html}.
\bibitem{method}
Brian Cameron. \emph{Operator Archaeology: A Protocol for Bounded
Reconstruction of Procedural Media}. Companion methodology manuscript, 2026.
\url{https://github.com/yippibrian/history/tree/main/00-methodology/01-operator-archaeology}.
\bibitem{brack}
Lis Brack-Bernsen. \emph{Babylonian astronomy: a new understanding of column
$\Phi$}. Archive for History of Exact Sciences 74 (2020).
\url{https://doi.org/10.1007/s00407-020-00254-z}.
\bibitem{mesopotamia}
Brian Cameron. \emph{Mesopotamian Astronomical Procedure}.
Companion historical manuscript, 2026.
\url{https://github.com/yippibrian/history/tree/main/01-ancient-operator-systems/13-mesopotamian-astronomy}.
\bibitem{urboard}
British Museum. \emph{Game-board}, museum number 1928,1009.378.
Collection record.
\url{https://www.britishmuseum.org/collection/object/W_1928-1009-378}.
\bibitem{urrules}
British Museum. \emph{Tablet}, museum number Rm.III.6.b.
Collection record for the explanatory diagram and rules of the twenty-square game.
\url{https://www.britishmuseum.org/collection/object/W_Rm-III-6-b}.
See also Irving L. Finkel, \emph{On the Rules for the Royal Game of Ur}, in
\emph{Ancient Board Games in Perspective}, edited by Finkel, British Museum
Press, 2007, pp. 16--32.
\bibitem{lucarelli}
Rita Lucarelli. \emph{The guardian-demons of the Book of the Dead}.
British Museum Studies in Ancient Egypt and Sudan 15 (2010): 85--102,
especially the discussion of spells 144--147 and their temple adaptations.
Accessible reproduction:
\url{https://www.readkong.com/page/the-guardian-demons-of-the-book-of-the-dead-2250019}.
\bibitem{freeth}
Tony Freeth et al. \emph{Decoding the ancient Greek astronomical calculator known
as the Antikythera Mechanism}. Nature 444 (2006): 587--591.
\url{https://doi.org/10.1038/nature05357}.
\bibitem{priani}
Ernesto Priani. \emph{Ramon Llull}. Stanford Encyclopedia of Philosophy.
Revised February 24, 2025; especially sections 1 and 5.
\url{https://plato.stanford.edu/entries/llull/}.
\bibitem{llull}
Ramon Llull. \emph{Ars brevis}, 1308. English translation by Yanis Dambergs,
December 2006, especially Part 2, Fourth Figure, and the definitions and rules.
Digital translation:
\url{https://www.labirintoermetico.com/12ArsCombinatoria/Llull_R_Ars_Brevis_(tr._inglese).pdf}.
Latin text, Archivio della Latinit\`a Italiana del Medioevo:
\url{https://alim.unisi.it/dl/resource/181}.
\bibitem{reeds}
Jim Reeds. \emph{Solved: The Ciphers in Book III of Trithemius's Steganographia}.
Cryptologia 22, no. 4 (1998): 291--317.
\url{https://doi.org/10.1080/0161-119891886948}.
\bibitem{trithemiusPaper}
Brian Cameron. \emph{Reconstructing the Operational Machine of Trithemius's
Steganographia}. Companion historical manuscript, 2026.
\url{https://github.com/yippibrian/history/tree/main/04-early-modern-cryptography/01-trithemius-steganographia}.
\bibitem{nlw}
National Library of Wales. \emph{Steganographia}, Peniarth MS 423D.
Digital exhibition: date, attribution, earlier copying history, and limits
of claims about Dee's use.
\url{https://www.library.wales/discover-learn/digital-exhibitions/manuscripts/early-modern-period/steganographia}.
\bibitem{deeLayout}
John Dee. \emph{A True \& Faithful Relation of What Passed for Many Yeers
Between Dr. John Dee and Some Spirits}. Edited by Meric Casaubon, 1659.
Part I, entries for June 18--19, 1583, printed pp. 18--21, especially 20--21.
Digital transcription hosted by Joseph H. Peterson:
\url{https://www.esotericarchives.com/dee/tfr/tfr1.htm}.
\bibitem{heptarchia}
John Dee. \emph{De Heptarchia Mystica}. Digital edition by Joseph H. Peterson,
based on British Library Sloane MS 3191; chapter 3, fols. 38r--38v,
April 28--29, 1583 arrangement.
\url{https://www.esotericarchives.com/dee/hm.htm}.
\bibitem{deePaper}
Brian Cameron. \emph{The Angelic Tables of John Dee}.
Companion historical manuscript, 2026.
\url{https://github.com/yippibrian/history/tree/main/04-early-modern-cryptography/02-john-dee}.
\bibitem{leibniz}
Brandon C. Look. \emph{Gottfried Wilhelm Leibniz}.
Stanford Encyclopedia of Philosophy, especially the intellectual
biography and discussion of the universal characteristic.
\url{https://plato.stanford.edu/entries/leibniz/}.
\bibitem{vaswani}
Ashish Vaswani et al. \emph{Attention Is All You Need}.
Advances in Neural Information Processing Systems 30 (2017).
\url{https://arxiv.org/abs/1706.03762}.
\bibitem{ouyang}
Long Ouyang et al. \emph{Training language models to follow instructions with
human feedback}. Advances in Neural Information Processing Systems
35 (2022): 27730--27744.
\url{https://arxiv.org/abs/2203.02155}.
\bibitem{reliability}
Brian Cameron. \emph{Structural Reliability Under Projection:
Representational Loss, Reconstructibility, and Admissibility}.
Working paper, 2026.
\url{https://github.com/yippibrian/reconstructibility-under-projection/tree/main/01-reconstruction/01-reconstruction-under-projection}.
\bibitem{rprp}
Brian Cameron. \emph{RPRP: Reconstructibility-Preserving Reasoning Under Projection}.
Working paper, 2026.
\url{https://github.com/yippibrian/reconstructibility-under-projection/tree/main/01-reconstruction/04-rprp}.
\bibitem{bundles}
Brian Cameron. \emph{Governed Reasoning Bundles}.
Working paper, 2026.
\url{https://github.com/yippibrian/reconstructibility-under-projection/tree/main/02-shared/05-governed-bundles}.
\bibitem{caves}
Brian Cameron. \emph{Between Caves}.
Working paper, 2026.
\url{https://github.com/yippibrian/reconstructibility-under-projection/tree/main/03-philosophy/01-caves}.
\bibitem{controller}
Brian Cameron. \emph{The Controller}.
Working paper, 2026.
\url{https://github.com/yippibrian/reconstructibility-under-projection/tree/main/03-philosophy/03-controller}.
\end{thebibliography}
\end{document}

